\section{Experimental campaign}
\label{sec:experiments}

\subsection{Execution environment}

The entire campaign runs as a matrix of jobs in a public continuous-integration
workflow on a single standard runner with four virtual CPUs, 16~GB of RAM and
local SSD storage. No dedicated hardware, cloud credits or multi-host cluster
are used. This is a deliberate design choice: it makes the campaign
reproducible by any reader at zero cost, and it exercises the \RNM{} protocol of
Section~\ref{sec:rnm} in exactly the regime it was designed for.

The consensus platform is Hyperledger Besu with QBFT, deployed as containers
from a generated genesis file. Each validator container is granted a fixed CPU
quota~$\q$ and memory limit, and the JVM heap is capped so that
$\n$~validators fit within the host memory. The CI matrix dimension maps
one-to-one onto the experimental design matrix, and every job uploads its raw
measurements as an artefact; a final aggregation job assembles the dataset.

\subsection{Load generation}

Client concurrency is realized by $\cc$ independent asynchronous senders, each
owning a distinct externally-owned account and maintaining its own nonce pool.
Per-sender nonce ownership removes the single-threaded submission bottleneck
that limits naive load generators and would otherwise saturate the measurement
before the consensus layer is stressed. Throughput is computed from on-chain
block contents rather than from client acknowledgements, so that submission
failures cannot inflate the reported value. The first and last windows of each
run are discarded as warm-up and drain.

\subsection{Fault injection and detection}

\paragraph{Injection.}
Faulty behaviour follows Definition~\ref{def:faults} and is injected at the
network layer with kernel traffic-control queueing disciplines applied to the
selected containers: a loss probability implements omission, and a duplication
probability implements message repetition. Because the peer-to-peer transport is
encrypted, injection cannot be made selective with respect to message type; we
therefore claim only the omission/duplication fault class and state this
limitation explicitly in Section~\ref{sec:limitations}.

\paragraph{Detection.}
The monitor polls the platform's signer-metrics interface, which reports the
number of blocks proposed by each validator over a range of block heights. A
validator whose participation over the window falls below a threshold is
flagged. Because the set of degraded containers is known by construction, every
run yields ground-truth labels, from which we estimate the detection
probability~$\pd$, the false-positive rate~$\pf$, the receiver operating
characteristic, and the exchangeable correlation~$\rho$ of
Assumption~\ref{as:rho}. This procedure requires no modification of the
consensus software.

\subsection{Design matrix}

Table~\ref{tab:design} lists the experiments. Each cell is replicated three
times and each run measures a fixed steady-state interval after warm-up.

\begin{table}[!t]
\caption{Experimental design. $\q$ is the per-validator CPU quota in cores,
$\cc$ the number of independent senders, RTT the emulated round-trip latency,
and $\ell$ the injected loss/duplication level}
\label{tab:design}
\centering
\small
\begin{tabular}{llp{74mm}}
\headline
Id & Purpose & Factor levels \\
\headline
X1 & Consensus exponent $\beta$
   & $\n\in\{4,5,7,9,11,13,16\}$; $\q$ fixed; $\cc$ fixed \\
X2 & $\q$-invariance (Statement~\ref{st:qinvariance})
   & $\n\in\{4,7,13\}\times\q\in\{0.20,0.25,0.33\}$ \\
X3 & Client factor $\gamma,\csat,\theta$
   & $\cc\in\{1,2,4,8,16,32,64\}$; $\n$ fixed \\
X4 & Latency sensitivity $\beta(\text{RTT})$
   & RTT $\in\{0,25,50,100,200\}$~ms $\times\ \n\in\{7,13\}$ \\
X5 & Detector characterization $\pd,\pf,\rho$
   & $\ell\in\{0,10,20,25\}\%$ $\times\ \n\in\{7,13\}$ \\
X6 & Extrapolation to large $\n$
   & discrete-event simulation calibrated on X1--X4 \\
X7 & Coverage validation
   & calibrate on $\n\le9$, predict held-out $\n\in\{13,16\}$ \\
X8 & Sizing ablation
   & linear, single-factor~\cite{peniak2026a}, USL~\cite{gunther2007}, ours \\
X9 & Cross-implementation replication
   & X1 and X3 repeated on a second QBFT-family client \\
\headline
\end{tabular}
\end{table}

\subsection{Confounding test}

Theorem~\ref{thm:identifiability} is tested directly rather than argued. From
the factorial grid of X1 and X3 we estimate $\beta$ and $\gamma$ with
orthogonal factors. We then \emph{re-sample} the same measurements along an
artificial concurrency path $\log\cc=\lambda\log\n+\mu$ for several values of
$\lambda$, fit the single-factor model to each resampled subset, and compare the
resulting $\ahat$ against the prediction $\gamma\lambda-\beta$
of~\eqref{eq:confounding}. Agreement over a range of $\lambda$, including a
value at which $\ahat$ changes sign, constitutes an empirical confirmation of
the theorem.

Experiment X9 repeats the exercise on a second client of the same protocol
family. The coefficients $(\beta,\gamma)$ are expected to differ; the line
$\ahat=\gamma\lambda-\beta$ is not, since the theorem concerns the estimator and
not the system. This separates a claim about least squares from a claim about
Hyperledger Besu, and it is the cheapest available check that our conclusions
are not an artefact of one implementation.

\subsection{Reproducibility}

The workflow definition, network generator, load generator, fault-injection
scripts, detector, analysis code and raw measurement artefacts are published
together with the paper. All numeric values quoted in
Section~\ref{sec:results} are emitted by the analysis pipeline into a generated
file that the manuscript includes, so that no reported number is transcribed by
hand.
